% common-scrbook/preamble/setup.tex
%
% Base preamble for the scrbook target. Loaded by main.tex after
% \documentclass{scrbook}. Establishes language, math, fonts, hyperlinks,
% cross-references, and color palette. Segment-specific commands
% (\segmenthead, \eqtag, \segmentfield) and environments (workingnotes,
% callout) live in segment.tex / environments.tex.

%----------------------------------------------------------------------------
% Language & quoting
%----------------------------------------------------------------------------
\usepackage[english]{babel}
\usepackage[english=american]{csquotes}

%----------------------------------------------------------------------------
% Math — amsmath family, mathtools (which patches several amsmath rough
% edges), then the unicode-math layer for fontspec-driven math fonts.
%----------------------------------------------------------------------------
\usepackage{amsmath}
\usepackage{amssymb}
\usepackage{mathtools}
\usepackage{amsthm}
\usepackage{unicode-math}

%----------------------------------------------------------------------------
% Fonts.
%
% Default is STIX Two — a single coordinated text + math family that ships
% with current TeX Live and renders with the visual register of a classical
% mathematical monograph (the AMS GSM series, Folland's Real Analysis,
% Lang's Algebra). Both faces are open-source and machine-portable.
%
% To swap fonts: replace the three blocks below. fontspec finds OS fonts via
% fontconfig and TeX Live fonts via its texmf tree; either source works.
%----------------------------------------------------------------------------
% ── Body serif: Adobe Garamond Premier Pro + Joseph's modified Alt ──
%
% Joseph licensed the Adobe Garamond Premier Pro family (~36 OTFs covering
% the optical-size design — caption / text / subhead / display — across
% upright/italic/medium/semibold/bold). He then made GaramondPremrProAlt
% (a FontForge-edited version of the text-optical-size regular weight)
% with revised dash, question mark, and tighter kerning around quote /
% punctuation pairs (e.g. .", ?', etc.). Source files (*.sfd) live under
% ~/src/_ref/axiomata/ and ~/src/_ref/backup-simplex-stuff/.
%
% Setup convention:
%   - Body upright: GaramondPremrProAlt (Joseph's modifications)
%   - Italic / Bold / Bold-Italic: Adobe stock garamondpremrpro-{it,bd,bdit}
%   - Optical sizing via SizeFeatures: caption (≤ 8.9pt), text (8.9-14.9pt,
%     designed at 11pt), subhead (14.9-22.9pt), display (≥ 22.9pt)
%   - HarfBuzz renderer for thorough OpenType + kerning resolution
%   - lining figures + oldstyle numbers per gpp-norm intent
%
% Path is repo-root-relative; bin/build-monograph runs lualatex with
% chdir=ROOT_DIR so this resolves to mono/scrbook/fonts/garamond/ in the
% checked-out source tree. Font files are committed under
% mono/scrbook/fonts/garamond/.
%
% Fallback chain for missing-glyph coverage MUST be registered BEFORE
% \setmainfont references it via RawFeature={fallback=...}. luaotfload's
% add_fallback registers a named font list; fontspec consults the chain
% when the primary font lacks a glyph. Order matters: search the most-
% aesthetically-coordinated fallback first (Noto Serif pairs with Garamond),
% then broader-coverage fonts (FreeSerif covers most BMP + supplemental
% planes; DejaVu Sans is the catch-all for symbols / extended Latin).
\directlua{
  if luaotfload and luaotfload.add_fallback then
    luaotfload.add_fallback('serif-fallback', {
      'NotoSerif:mode=harf;script=latn;language=dflt;',
      'FreeSerif:mode=harf;script=latn;language=dflt;',
      'DejaVuSans:mode=harf;script=latn;language=dflt;',
    })
  end
}

\setmainfont{GaramondPremrProAlt}[
  Path               = ./mono/scrbook/fonts/garamond/,
  Extension          = .otf,
  Renderer           = HarfBuzz,
  Ligatures          = {TeX, Common},
  Numbers            = {Proportional, Lining},  % uppercase / all same height —
                                                 % more authoritative for
                                                 % technical / numbered
                                                 % content (table numbers,
                                                 % equation numbers, page
                                                 % numbers, dates in prose).
                                                 % Tabular spacing for
                                                 % numerical table cells is
                                                 % a deferred separate
                                                 % capability — see TODO in
                                                 % bin/lib/typeset_scrbook.rb
                                                 % convert_table.
  RawFeature         = {fallback=serif-fallback},
  ItalicFont         = garamondpremrpro-it,
  BoldFont           = garamondpremrpro-bd,
  BoldItalicFont     = garamondpremrpro-bdit,
  % Optical-size dispatch — fontspec swaps the file based on requested
  % point size. Ranges follow Adobe's design-intent boundaries (preserved
  % from Joseph's original tuftish-fonts setup). The body range is the
  % default (no override); other ranges swap in the corresponding optical
  % variants for upright + italic + bold + bold-italic.
  SizeFeatures = {
    {Size = -8.9,
     Font           = garamondpremrpro-capt,
     ItalicFont     = garamondpremrpro-itcapt,
     BoldFont       = garamondpremrpro-bdcapt,
     BoldItalicFont = garamondpremrpro-bditcapt},
    {Size = 8.9-14.9},  % body — defaults from outer setmainfont
    {Size = 14.9-22.9,
     Font           = garamondpremrpro-subh,
     ItalicFont     = garamondpremrpro-itsubh,
     BoldFont       = garamondpremrpro-bdsubh,
     BoldItalicFont = garamondpremrpro-bditsubh},
    {Size = 22.9-,
     Font           = garamondpremrpro-disp,
     ItalicFont     = garamondpremrpro-itdisp,
     BoldFont       = garamondpremrpro-bddisp,
     BoldItalicFont = garamondpremrpro-bditdisp},
  },
]
% Math font — STIX Two Math. Garamond Premier Pro has no coordinated math
% companion; STIX Two Math renders Greek + symbols cleanly enough alongside
% most modern serifs and is the standard choice for math monographs whose
% body family lacks a math sibling.
\setmathfont{STIX Two Math}

% Sans-serif: Inter — used by KOMA scrbook for chapter / section /
% subsection headings. Variable font; same recipe as Lora — explicit
% file-path + axis instances per style, HarfBuzz renderer for the
% +axis feature to take effect.
% Inter's filenames embed a literal comma (Inter-VariableFont_opsz,wght.ttf).
% fontspec's option parser splits values on commas, so comma-bearing values
% MUST be brace-quoted ({...}); backslash-escaping (\,) is interpreted as
% the thinspace command and corrupts the filename — which then leaks
% fontspec's internal identifier strings as visible content on the title
% page when font resolution fires (the Inter \maketitle / \publishers
% rendering was the canary).
\setsansfont{Inter-VariableFont_opsz,wght}[
  Path               = /Library/Fonts/,
  Extension          = .ttf,
  Renderer           = HarfBuzz,
  Ligatures          = TeX,
  Scale              = MatchLowercase,
  ItalicFont         = {Inter-Italic-VariableFont_opsz,wght},
  BoldFont           = {Inter-VariableFont_opsz,wght},
  BoldItalicFont     = {Inter-Italic-VariableFont_opsz,wght},
  UprightFeatures    = {RawFeature={+axis={wght=400}}},
  BoldFeatures       = {RawFeature={+axis={wght=700}}},
  ItalicFeatures     = {RawFeature={+axis={wght=400}}},
  BoldItalicFeatures = {RawFeature={+axis={wght=700}}},
]

% Monospace (technically proportional): FiraCode Nerd Font Propo — variable-
% width variant with Nerd Font icon glyphs preserved. Inline \texttt{} content
% (slug names, file paths, identifiers) reads with FiraCode's distinctive
% character forms without the cramped fixed-cell feel of true monospace,
% while keeping the code-like aesthetic register. Locally installed via
% ~/Library/Fonts/ — fontspec finds it through fontconfig.
\setmonofont{FiraCode Nerd Font}[
  Scale       = MatchLowercase,
  LetterSpace = -5,           % tighten the inter-character gap; FiraCode reads
                              % loose by default at small sizes. Negative values
                              % pull characters closer; tune by feel.
]

%----------------------------------------------------------------------------
% GFM/Obsidian math-compat macros — same shims as the kaobook target.
% Source segments use \lt and \gt because raw < and > break GitHub's math
% renderer (interpreted as HTML tags). LaTeX math knows neither — define
% them as the obvious operators.
%----------------------------------------------------------------------------
\providecommand{\lt}{<}
\providecommand{\gt}{>}

%----------------------------------------------------------------------------
% Microtypography
%----------------------------------------------------------------------------
\usepackage{microtype}

%----------------------------------------------------------------------------
% Widow / orphan suppression. Same discipline as the kaobook target — a
% stranded single line at top or bottom of a page reads as broken; the
% occasional extra whitespace is the right trade.
%----------------------------------------------------------------------------
\widowpenalty=10000
\clubpenalty=10000

%----------------------------------------------------------------------------
% String predicates — used by segment.tex for the eq-tag dispatch and any
% conditional rendering.
%----------------------------------------------------------------------------
\usepackage{xstring}

%----------------------------------------------------------------------------
% \AlphAlph — \Alph with overflow protection. Standard \Alph errors at
% values > 26 (Latin alphabet only). \AlphAlph gives A–Z for 1–26, then
% AA–AZ, BA–BZ, … for 27 onward. Used by appendix-chapter numbering, where
% AAT's Details appendix alone has more than 26 chapter-level entries — so
% the native scrbook \appendix \Alph numbering hits "Counter too large" at
% the 27th appendix and corrupts every appendix and ToC entry past Z.
%
% Deliberately byte-identical to the kaobook target's definition
% (mono/kaobook/preamble/setup.tex). The two preamble trees are separate by
% design; if you change one, change both — they encode the same scheme so
% the target-agnostic assembled markdown renders the same appendix labels
% under either target.
%
% Formula: for n in [27, 26+26²], chunk = (n-1) ÷ 26, rem = ((n-1) mod 26)+1,
% emitted as \@Alph{chunk}\@Alph{rem}. \numexpr handles the integer math.
%----------------------------------------------------------------------------
\makeatletter
\newcommand{\AlphAlph}[1]{%
  \ifnum#1<27%
    \@Alph{#1}%
  \else
    \@Alph{\numexpr(#1-1)/26\relax}\@Alph{\numexpr#1-26*((#1-1)/26)\relax}%
  \fi
}
\makeatother

%----------------------------------------------------------------------------
% Color palette. Modest by design — link/cross-ref navy, a quiet rule-gray,
% a muted sand for the working-notes register. No status-tier colors:
% scrbook target drops the status badges entirely.
%----------------------------------------------------------------------------
\usepackage{xcolor}
\definecolor{rule-gray}      {HTML}{B8B8B8}
\definecolor{link-navy}      {HTML}{1F3A6F}
\definecolor{wn-rule}        {HTML}{8C8676}
\definecolor{wn-fg}          {HTML}{6E6354}
\definecolor{callout-rule}   {HTML}{6B7B95}

%----------------------------------------------------------------------------
% pdfpages — for the cover page (rsvg-convert produces a single-page PDF
% that we \includepdf at the very front).
%----------------------------------------------------------------------------
\usepackage{pdfpages}

%----------------------------------------------------------------------------
% needspace — \needspace{N\baselineskip} prevents a heading from being
% orphaned at page bottom by forcing a page break first when N lines aren't
% available below.
%----------------------------------------------------------------------------
\usepackage{needspace}

%----------------------------------------------------------------------------
% Tables — same toolkit as the kaobook target so AsfLatex's table converter
% (in segment_renderer.rb) renders identically here.
%----------------------------------------------------------------------------
\usepackage{tabularx}
\usepackage{xltabular}    % longtable + tabularx — page-breakable, width-aware.
                          % Used by AsfScrbookLatex's convert_table override
                          % (the kaobook target sticks with plain tabularx
                          % because xltabular collides with kaobook's
                          % margin-note machinery, see segment_renderer.rb).
\usepackage{booktabs}
\usepackage{array}
\usepackage{caption}
\captionsetup{justification=raggedright,singlelinecheck=false}
% Caption registers under the table (LongTable convention via \caption inside
% the env still places it where \caption is called). We emit \caption as the
% first line of xltabular content, so the caption appears ABOVE the table
% header — consistent with the kaobook target's \captionof{table}{...} above.

% Figure embeds (segment_renderer.rb convert_img). Markdown figure
% embeds resolve to an in-flow \captionof{figure} atom whose graphic
% is a PDF: a TikZ figure's `.tex` (source of truth) is recompiled to
% a gitignored `.mono.pdf` by the resolver when it changes, then
% \includegraphics'd. Only graphicx is needed host-side (no
% \includestandalone — structurally incompatible with this pipeline;
% see msc/figure-pipeline-buildout-2026-05-18.md).
\usepackage{graphicx}

%----------------------------------------------------------------------------
% Equation + table numbering reset per chapter — same convention as the
% kaobook target (3.5, 3.6, …, then 4.1 in next chapter).
%----------------------------------------------------------------------------
\numberwithin{equation}{chapter}
\numberwithin{table}{chapter}
% Figures: FLAT volume-wide numbering (Figure 1, 2, …), deliberately
% NOT per-chapter like equations/tables. The Introduction is an
% unnumbered \addchap, so a \thechapter-prefixed figure number
% collapses to "Figure :" there; figures are also far sparser than
% equations/tables, so a flat scheme reads cleanly and is front-
% matter-safe. This is the figure-numbering spec decision Joseph
% delegated ("decide + build it in") — flagged in
% msc/figure-pipeline-buildout-2026-05-18.md for his review; trivially
% switchable to \numberwithin{figure}{chapter} if he prefers per-chapter
% and accepts special-casing the Introduction. \counterwithout removes
% the book-class per-chapter RESET (else each chapter restarts at
% "Figure 1"); the renew drops the \thechapter prefix.
\usepackage{chngcntr}
\counterwithout{figure}{chapter}
\renewcommand{\thefigure}{\arabic{figure}}

%----------------------------------------------------------------------------
% Hyperlinks + cross-references. hyperref must load late (it patches many
% other packages). cleveref must load AFTER hyperref so its \cref renders
% as a clickable link.
%----------------------------------------------------------------------------
\usepackage[
  colorlinks = true,
  linkcolor  = link-navy,
  citecolor  = link-navy,
  urlcolor   = link-navy,
  filecolor  = link-navy,
  pdfborder  = {0 0 0},
]{hyperref}

\usepackage[capitalize,noabbrev]{cleveref}

% Cross-ref display: bare number for sections and chapters. Resolved cross-
% refs in the assembled markdown look like `[Definition 1.4](#def-foo)`;
% the converter emits `\cref{seg:def-foo}`. We want LaTeX to produce the
% number "1.4" — the type word ("Definition") was kept in the markdown's
% visible link text and survives in source prose like "see Definition
% \cref{seg:def-foo}".
%
% IMPORTANT: cleveref's format placeholder for the stored label value is
% `#1`, NOT `\thesection`. `#2#1#3` is the canonical
% "hyperref-open + value + hyperref-close" idiom.
\AtBeginDocument{%
  \crefformat{section}{#2#1#3}%
  \Crefformat{section}{#2#1#3}%
  \crefrangeformat{section}{#3#1#4--#5#2#6}%
  \Crefrangeformat{section}{#3#1#4--#5#2#6}%
  \crefmultiformat{section}{#2#1#3}{, #2#1#3}{, #2#1#3}{, #2#1#3}%
  \Crefmultiformat{section}{#2#1#3}{, #2#1#3}{, #2#1#3}{, #2#1#3}%
  % Appendix segments are chapters; mirror the section convention so
  % chapter and section cross-refs read identically in running text.
  \crefformat{chapter}{#2#1#3}%
  \Crefformat{chapter}{#2#1#3}%
  \crefrangeformat{chapter}{#3#1#4--#5#2#6}%
  \Crefrangeformat{chapter}{#3#1#4--#5#2#6}%
  \crefmultiformat{chapter}{#2#1#3}{, #2#1#3}{, #2#1#3}{, #2#1#3}%
  \Crefmultiformat{chapter}{#2#1#3}{, #2#1#3}{, #2#1#3}{, #2#1#3}%
}

%----------------------------------------------------------------------------
% \ifdefempty — used by main.tex to detect empty cover-path. Provided by
% etoolbox, which KOMA pulls in transitively, but load explicitly so the
% conditional always resolves.
%----------------------------------------------------------------------------
\usepackage{etoolbox}
